\documentclass[tikz,border=8pt]{standalone}
\usepackage{tikz}
\usepackage{amsmath}
\usepackage{amssymb}
\usepackage{pgfplots}
\pgfplotsset{compat=1.18}
\usepackage{ctex}
\usetikzlibrary{calc,positioning,arrows.meta,trees,matrix}

% 归并排序分治递归树（LC 148 Sort List）
% 数组 [3,1,4,1,5,9,2,6]
% 白色节点：分（divide），绿色节点：合并结果

\begin{document}
\begin{tikzpicture}[
  every node/.style={font=\small},
  divide node/.style={draw,rectangle,rounded corners=2pt,minimum width=2.8cm,minimum height=0.55cm,
                       fill=white,draw=gray!70,thick},
  merge node/.style={draw,rectangle,rounded corners=2pt,minimum width=2.8cm,minimum height=0.55cm,
                      fill=green!20,draw=green!60!black,thick},
  leaf node/.style={draw,rectangle,rounded corners=2pt,minimum width=0.8cm,minimum height=0.55cm,
                     fill=blue!10,draw=blue!50,thick},
  arr/.style={->,>=Stealth,thick},
  back arr/.style={->,>=Stealth,thick,green!60!black},
]

%% ── 标题 ────────────────────────────────────────────────────
\node[font=\large\bfseries] at (0, 1.8) {归并排序分治树（LC 148 Sort List）};
\node[font=\small,gray] at (0, 1.15)
  {数组 $[3,1,4,1,5,9,2,6]$，白色=分 / 绿色=合并结果};

%% ════════════════════════════════════════════════════════════
%% 层 0：完整数组（白色，分）
%% ════════════════════════════════════════════════════════════
\node[divide node,minimum width=5.5cm] (L0) at (0, 0) {3, 1, 4, 1, 5, 9, 2, 6};

%% ════════════════════════════════════════════════════════════
%% 层 1：两个半段（白色，分）
%% ════════════════════════════════════════════════════════════
\node[divide node,minimum width=2.6cm] (L1a) at (-4.5, -1.8) {3, 1, 4, 1};
\node[divide node,minimum width=2.6cm] (L1b) at ( 4.5, -1.8) {5, 9, 2, 6};

\draw[arr] (L0.south) -- ++(0,-0.3) -| (L1a.north);
\draw[arr] (L0.south) -- ++(0,-0.3) -| (L1b.north);

%% ════════════════════════════════════════════════════════════
%% 层 2：四段（白色，分）
%% ════════════════════════════════════════════════════════════
\node[divide node,minimum width=1.5cm] (L2a) at (-7.0, -3.6) {3, 1};
\node[divide node,minimum width=1.5cm] (L2b) at (-2.0, -3.6) {4, 1};
\node[divide node,minimum width=1.5cm] (L2c) at ( 2.0, -3.6) {5, 9};
\node[divide node,minimum width=1.5cm] (L2d) at ( 7.0, -3.6) {2, 6};

\draw[arr] (L1a.south) -- ++(0,-0.3) -| (L2a.north);
\draw[arr] (L1a.south) -- ++(0,-0.3) -| (L2b.north);
\draw[arr] (L1b.south) -- ++(0,-0.3) -| (L2c.north);
\draw[arr] (L1b.south) -- ++(0,-0.3) -| (L2d.north);

%% ════════════════════════════════════════════════════════════
%% 层 3：单元素（蓝色叶子）
%% ════════════════════════════════════════════════════════════
\node[leaf node] (L3a) at (-8.2, -5.4) {3};
\node[leaf node] (L3b) at (-5.8, -5.4) {1};
\node[leaf node] (L3c) at (-3.2, -5.4) {4};
\node[leaf node] (L3d) at (-0.8, -5.4) {1};
\node[leaf node] (L3e) at ( 0.8, -5.4) {5};
\node[leaf node] (L3f) at ( 3.2, -5.4) {9};
\node[leaf node] (L3g) at ( 5.8, -5.4) {2};
\node[leaf node] (L3h) at ( 8.2, -5.4) {6};

\draw[arr] (L2a.south) -- ++(0,-0.3) -| (L3a.north);
\draw[arr] (L2a.south) -- ++(0,-0.3) -| (L3b.north);
\draw[arr] (L2b.south) -- ++(0,-0.3) -| (L3c.north);
\draw[arr] (L2b.south) -- ++(0,-0.3) -| (L3d.north);
\draw[arr] (L2c.south) -- ++(0,-0.3) -| (L3e.north);
\draw[arr] (L2c.south) -- ++(0,-0.3) -| (L3f.north);
\draw[arr] (L2d.south) -- ++(0,-0.3) -| (L3g.north);
\draw[arr] (L2d.south) -- ++(0,-0.3) -| (L3h.north);

%% ════════════════════════════════════════════════════════════
%% 层 4（合并 2 元素）：绿色节点，含排序结果
%% ════════════════════════════════════════════════════════════
\node[merge node,minimum width=1.5cm] (M2a) at (-7.0, -7.2) {1, 3};
\node[merge node,minimum width=1.5cm] (M2b) at (-2.0, -7.2) {1, 4};
\node[merge node,minimum width=1.5cm] (M2c) at ( 2.0, -7.2) {5, 9};
\node[merge node,minimum width=1.5cm] (M2d) at ( 7.0, -7.2) {2, 6};

\draw[back arr] (L3a.south) -- ++(0,-0.3) -| (M2a.north);
\draw[back arr] (L3b.south) -- ++(0,-0.3) -| (M2a.north);
\draw[back arr] (L3c.south) -- ++(0,-0.3) -| (M2b.north);
\draw[back arr] (L3d.south) -- ++(0,-0.3) -| (M2b.north);
\draw[back arr] (L3e.south) -- ++(0,-0.3) -| (M2c.north);
\draw[back arr] (L3f.south) -- ++(0,-0.3) -| (M2c.north);
\draw[back arr] (L3g.south) -- ++(0,-0.3) -| (M2d.north);
\draw[back arr] (L3h.south) -- ++(0,-0.3) -| (M2d.north);

% 合并标注
\node[font=\tiny,green!60!black] at (-7.0,-6.6) {merge};
\node[font=\tiny,green!60!black] at (-2.0,-6.6) {merge};
\node[font=\tiny,green!60!black] at ( 2.0,-6.6) {merge};
\node[font=\tiny,green!60!black] at ( 7.0,-6.6) {merge};

%% ════════════════════════════════════════════════════════════
%% 层 5（合并 4 元素）：绿色节点
%% ════════════════════════════════════════════════════════════
\node[merge node,minimum width=2.6cm] (M1a) at (-4.5, -9.0) {1, 1, 3, 4};
\node[merge node,minimum width=2.6cm] (M1b) at ( 4.5, -9.0) {2, 5, 6, 9};

\draw[back arr] (M2a.south) -- ++(0,-0.3) -| (M1a.north);
\draw[back arr] (M2b.south) -- ++(0,-0.3) -| (M1a.north);
\draw[back arr] (M2c.south) -- ++(0,-0.3) -| (M1b.north);
\draw[back arr] (M2d.south) -- ++(0,-0.3) -| (M1b.north);

\node[font=\tiny,green!60!black] at (-4.5,-8.4) {merge};
\node[font=\tiny,green!60!black] at ( 4.5,-8.4) {merge};

%% ════════════════════════════════════════════════════════════
%% 层 6（最终合并 8 元素）：绿色节点，最终结果
%% ════════════════════════════════════════════════════════════
\node[merge node,minimum width=5.5cm,draw=green!70!black,very thick,fill=green!30]
  (M0) at (0, -10.8) {1, 1, 2, 3, 4, 5, 6, 9};

\draw[back arr,very thick] (M1a.south) -- ++(0,-0.3) -| (M0.north);
\draw[back arr,very thick] (M1b.south) -- ++(0,-0.3) -| (M0.north);

\node[font=\small,green!60!black,above=2pt of M0] {\textbf{最终 merge}};

%% ── 层级标注（左侧） ──────────────────────────────────────────
\node[font=\scriptsize,gray,align=right] at (-11.5,  0.0)  {分 Divide\\第 1 层};
\node[font=\scriptsize,gray,align=right] at (-11.5, -1.8)  {分 Divide\\第 2 层};
\node[font=\scriptsize,gray,align=right] at (-11.5, -3.6)  {分 Divide\\第 3 层};
\node[font=\scriptsize,gray,align=right] at (-11.5, -5.4)  {叶 Leaf\\（单元素）};
\node[font=\scriptsize,green!60!black,align=right] at (-11.5, -7.2)  {合 Merge\\第 1 轮};
\node[font=\scriptsize,green!60!black,align=right] at (-11.5, -9.0)  {合 Merge\\第 2 轮};
\node[font=\scriptsize,green!60!black,align=right] at (-11.5,-10.8)  {合 Merge\\第 3 轮};

%% ── 说明框 ──────────────────────────────────────────────────
\node[draw=gray!50,fill=white,rounded corners=3pt,font=\small,inner sep=6pt,align=left]
  at (0,-12.4)
  {\textbf{复杂度：}时间 $O(n \log n)$，空间 $O(\log n)$（链表原地合并）\qquad
   \textbf{分治策略：}二分 → 递归排序 → 双指针合并};

\end{tikzpicture}
\end{document}
